% GENERATED FILE. Edit canonical lessons, then run npm run generate:books.
% canonical-source-hash: fnv1a64:74ed22c31b42a7ce
% canonical-lessons: KA-C309-hode, KA-C309-oduttidde, KA-C309-oduvaga, KA-R309-the-past-three-ways

\chapter{Did, Was Doing, Used To}
\label{ch:ka-did-was-doing-used-to}
\hlchaptermodality{309}

\begin{chapteropening}
\textbf{By the end of this chapter:} \emph{I can say what happened in the past for every person (hōde, bandaru), set the scene with ōduttidde and ōduvāga, say what used to happen, and say what there was with ittu.}

Complete the last lesson of chapter 309: tell what happened, what was going on and what used to happen, with the past by person, -uttidde, -uvāga and ittu.
\end{chapteropening}

\section[\texorpdfstring{\kn{ಹೋದೆ} (hōde)}{ಹೋದೆ (hōde)}]{\kn{ಹೋದೆ} (hōde) — the past, person by person}

\label{lesson:KA-C309-hode}



\begin{quote}
Before the new one: say the last two things you named, chalk and a ruler. (\emph{sīmesuṇṇa}, \emph{skēlu}.) Then \textquotedblleft{}I ate\textquotedblright{}, with the past bead you met in \textbf{\kn{ತಿಂದೆ}}. (\emph{tinde}.)

A story needs that past for everyone in it, not only for \textquotedblleft{}I\textquotedblright{}.
\end{quote}

\subsection*{You'll want to know: the past, person by person}
\begin{quote}
\textbf{\kn{ನಿನ್ನೆ} \kn{ನಾನು} \kn{ಊರಿಗೆ} \kn{ಹೋದೆ}.}
\end{quote}

\begin{quote}
\emph{ninne nānu ūrige hōde.} — \textquotedblleft{}Yesterday I went to my home town.\textquotedblright{}
\end{quote}

\textbf{\kn{ಹೋದೆ}} is \textquotedblleft{}I went\textquotedblright{}. Like \textbf{\kn{ತಿಂದೆ}}, it is three pieces: the stem, the past bead \textbf{-\kn{ದ}-}, and the person ending \textbf{-\kn{ಎ}}. Most verbs put an \textbf{-\kn{ಇ}-} in front of the bead: \textbf{\kn{ಮಾಡು}} $\to$ \textbf{\kn{ಮಾಡಿದ}-}. Hold that still and change only the ending:

\begin{itemize}
\raggedright
  \item \textbf{\kn{ನಾನು} \kn{ಮಾಡಿದೆ}} (\emph{nānu māḍide}) — I did
  \item \textbf{\kn{ನೀನು} \kn{ಮಾಡಿದೆ}} (\emph{nīnu māḍide}) — you did, to a friend or a child: the same ending as \textquotedblleft{}I\textquotedblright{}
  \item \textbf{\kn{ಅವನು} \kn{ಮಾಡಿದನು}} (\emph{avanu māḍidanu}) — he did
  \item \textbf{\kn{ಅವಳು} \kn{ಮಾಡಿದಳು}} (\emph{avaḷu māḍidaḷu}) — she did
  \item \textbf{\kn{ನಾವು} \kn{ಮಾಡಿದೆವು}} (\emph{nāvu māḍidevu}) — we did
  \item \textbf{\kn{ನೀವು} \kn{ಮಾಡಿದಿರಿ}} (\emph{nīvu māḍidiri}) — you did, to several people or with respect
  \item \textbf{\kn{ಅವರು} \kn{ಮಾಡಿದರು}} (\emph{avaru māḍidaru}) — they did; also \textquotedblleft{}he\textquotedblright{} or \textquotedblleft{}she\textquotedblright{}, said with respect
\end{itemize}

\textquotedblleft{}It\textquotedblright{} steps out of line: \textbf{\kn{ಅದು} \kn{ಮಾಡಿತು}} (\emph{adu māḍitu}), with \textbf{-\kn{ಇತು}} and no bead. Two \textquotedblleft{}it\textquotedblright{} verbs you will hear every day: \textbf{\kn{ಆಯಿತು}} (\emph{āyitu}, \textquotedblleft{}it happened, it became\textquotedblright{}) and \textbf{\kn{ಬಂತು}} (\emph{bantu}, \textquotedblleft{}it came\textquotedblright{}).

\subsection*{You'll want to know: the stems that change}
A few of the commonest verbs change their stem for the past. They are worth learning as pairs:

\begin{itemize}
\raggedright
  \item \textbf{\kn{ಬರು}} $\to$ \textbf{\kn{ಬಂದೆ}} (\emph{bande}), \textquotedblleft{}I came\textquotedblright{} (\textbf{\kn{ಬರು}} is the shape of \textbf{\kn{ಬಾ}} that carries endings)
  \item \textbf{\kn{ಹೋಗು}} $\to$ \textbf{\kn{ಹೋದೆ}} (\emph{hōde}), \textquotedblleft{}I went\textquotedblright{}; it went: \textbf{\kn{ಹೋಯಿತು}} (\emph{hōyitu})
  \item \textbf{\kn{ಕೊಡು}} $\to$ \textbf{\kn{ಕೊಟ್ಟೆ}} (\emph{koṭṭe}), \textquotedblleft{}I gave\textquotedblright{}
  \item \textbf{\kn{ಇರು}} $\to$ \textbf{\kn{ಇದ್ದೆ}} (\emph{idde}), \textquotedblleft{}I was\textquotedblright{}; it was: \textbf{\kn{ಇತ್ತು}} (\emph{ittu})
  \item \textbf{\kn{ತಿನ್ನು}} $\to$ \textbf{\kn{ತಿಂದೆ}} (\emph{tinde}), \textquotedblleft{}I ate\textquotedblright{}
\end{itemize}

The rest keep the \textbf{-\kn{ಇ}-} of \textbf{\kn{ಮಾಡಿದೆ}}: \textbf{\kn{ನೋಡಿದೆ}} (\emph{nōḍide}), \textquotedblleft{}I saw\textquotedblright{}; \textbf{\kn{ಓದಿದೆ}} (\emph{ōdide}), \textquotedblleft{}I read\textquotedblright{}. A stem that already ends in \textbf{-\kn{ಇ}} takes the bead straight on: \textbf{\kn{ಕುಡಿ}} $\to$ \textbf{\kn{ಕುಡಿದೆ}} (\emph{kuḍide}), \textquotedblleft{}I drank\textquotedblright{}.

\subsection*{Your turn}
Say these aloud:

\begin{itemize}
\raggedright
  \item \textquotedblleft{}Yesterday I went to my home town\textquotedblright{} — \emph{ninne nānu ūrige hōde}
  \item \textquotedblleft{}She came in the evening\textquotedblright{} — \emph{avaḷu sañje bandaḷu}
  \item \textquotedblleft{}We drank coffee\textquotedblright{} — \emph{nāvu kāphi kuḍidevu}
  \item \textquotedblleft{}I read a book. Afterwards I had a meal.\textquotedblright{} — \emph{nānu ondu pustaka ōdide. āmēle ūṭa māḍide.}
\end{itemize}

\begin{tcolorbox}[breakable,colback=teal!4,colframe=teal!35!black,title={Before you move on}]
\textquotedblleft{}I went\textquotedblright{}? (\textbf{\kn{ಹೋದೆ}}.) \textquotedblleft{}He came\textquotedblright{}? (\textbf{\kn{ಬಂದನು}}.) \textquotedblleft{}She gave\textquotedblright{}? (\textbf{\kn{ಕೊಟ್ಟಳು}}.) \textquotedblleft{}It happened\textquotedblright{}? (\textbf{\kn{ಆಯಿತು}}.) Which person ends in \textbf{-\kn{ತು}}? (\textbf{It}: \textbf{\kn{ಮಾಡಿತು}}, \textbf{\kn{ಬಂತು}}, \textbf{\kn{ಆಯಿತು}}.)
\end{tcolorbox}
\section[\texorpdfstring{\kn{ಓದುತ್ತಿದ್ದೆ} (ōduttidde)}{ಓದುತ್ತಿದ್ದೆ (ōduttidde)}]{\kn{ಓದುತ್ತಿದ್ದೆ} (ōduttidde) — what was going on, and what used to happen}

\label{lesson:KA-C309-oduttidde}



\begin{quote}
Say \textquotedblleft{}I read a book\textquotedblright{}. (\emph{nānu ondu pustaka ōdide}.) Then \textquotedblleft{}I was\textquotedblright{}, on its own. (\emph{idde}.)

The first is an event: it happened and it was over. A story also needs what was going on at the time.
\end{quote}

\subsection*{You'll want to know: -\kn{ಉತ್ತಿದ್ದೆ}, 'I was doing'}
\begin{quote}
\textbf{\kn{ನಿನ್ನೆ} \kn{ರಾತ್ರಿ} \kn{ಒಂಬತ್ತು} \kn{ಗಂಟೆಗೆ} \kn{ನಾನು} \kn{ಪುಸ್ತಕ} \kn{ಓದುತ್ತಿದ್ದೆ}.}
\end{quote}

\begin{quote}
\emph{ninne rātri ombattu gaṇṭege nānu pustaka ōduttidde.} — \textquotedblleft{}At nine last night I was reading a book.\textquotedblright{}
\end{quote}

\textbf{\kn{ಓದುತ್ತಿದ್ದೆ}} is two pieces you know: \textbf{\kn{ಓದುತ್ತ}}, the \textquotedblleft{}reading\textquotedblright{} inside \textbf{\kn{ಓದುತ್ತೇನೆ}}, and \textbf{\kn{ಇದ್ದೆ}}, \textquotedblleft{}I was\textquotedblright{}: \textquotedblleft{}I was, reading\textquotedblright{}. The ending is the ending of \textbf{\kn{ಇದ್ದೆ}}, so it changes the same way:

\begin{itemize}
\raggedright
  \item \textbf{\kn{ನಾನು} \kn{ಓದುತ್ತಿದ್ದೆ}} (\emph{nānu ōduttidde}) — I was reading
  \item \textbf{\kn{ಅವನು} \kn{ಓದುತ್ತಿದ್ದನು}} (\emph{avanu ōduttiddanu}) — he was reading
  \item \textbf{\kn{ಅವಳು} \kn{ಓದುತ್ತಿದ್ದಳು}} (\emph{avaḷu ōduttiddaḷu}) — she was reading
  \item \textbf{\kn{ಅವರು} \kn{ಓದುತ್ತಿದ್ದರು}} (\emph{avaru ōduttiddaru}) — they were reading
\end{itemize}

For \textquotedblleft{}it\textquotedblright{} the ending is the \textbf{\kn{ಇತ್ತು}} of \textquotedblleft{}it was\textquotedblright{}: \textbf{\kn{ಮಳೆ} \kn{ಬರುತ್ತಿತ್ತು}} (\emph{maḷe baruttittu}), \textquotedblleft{}it was raining\textquotedblright{}. \textbf{\kn{ಮಳೆ}} (\emph{maḷe}, \textquotedblleft{}rain\textquotedblright{}) \textquotedblleft{}comes\textquotedblright{} in Kannada.

\begin{grammarlens}[title={Grammar Lens: the same form for a habit}]
The same form says \textquotedblleft{}used to\textquotedblright{}, for something that happened many times over:

\begin{quote}
\textbf{\kn{ನನ್ನ} \kn{ಅಜ್ಜಿ} \kn{ಪ್ರತಿದಿನ} \kn{ಸಂಜೆ} \kn{ಒಂದು} \kn{ಕಥೆ} \kn{ಹೇಳುತ್ತಿದ್ದರು}.}
\end{quote}

\begin{quote}
\emph{nanna ajji pratidina sañje ondu kathe hēḷuttiddaru.} — \textquotedblleft{}My grandmother used to tell a story every evening.\textquotedblright{}
\end{quote}

\textbf{\kn{ಹೇಳುತ್ತಿದ್ದರು}} has the \textbf{-\kn{ಅರು}} of \textquotedblleft{}they\textquotedblright{}, because a grandmother is spoken of with respect, the way \textbf{\kn{ಅವರು}} can mean \textquotedblleft{}he\textquotedblright{} or \textquotedblleft{}she\textquotedblright{}.

Which meaning you have comes from the time words. One point in time, \textbf{\kn{ರಾತ್ರಿ} \kn{ಒಂಬತ್ತು} \kn{ಗಂಟೆಗೆ}}, gives \textquotedblleft{}was doing\textquotedblright{}. A word of repeating, \textbf{\kn{ಪ್ರತಿದಿನ}}, gives \textquotedblleft{}used to\textquotedblright{}: \textbf{\kn{ನಾನು} \kn{ಪ್ರತಿದಿನ} \kn{ಶಾಲೆಗೆ} \kn{ಹೋಗುತ್ತಿದ್ದೆ}} (\emph{nānu pratidina śālege hōguttidde}, \textquotedblleft{}I used to go to school every day\textquotedblright{}).
\end{grammarlens}

\subsection*{Your turn}
Say these aloud:

\begin{itemize}
\raggedright
  \item \textquotedblleft{}At nine last night I was reading a book\textquotedblright{} — \emph{ninne rātri ombattu gaṇṭege nānu pustaka ōduttidde}
  \item \textquotedblleft{}My grandmother used to tell a story every evening\textquotedblright{} — \emph{nanna ajji pratidina sañje ondu kathe hēḷuttiddaru}
  \item \textquotedblleft{}I used to go to school every day\textquotedblright{} — \emph{nānu pratidina śālege hōguttidde}
\end{itemize}

\begin{tcolorbox}[breakable,colback=teal!4,colframe=teal!35!black,title={Before you move on}]
\textquotedblleft{}I was reading\textquotedblright{}? (\textbf{\kn{ಓದುತ್ತಿದ್ದೆ}}.) \textquotedblleft{}It was raining\textquotedblright{}? (\textbf{\kn{ಮಳೆ} \kn{ಬರುತ್ತಿತ್ತು}}.) Which two pieces make \textbf{\kn{ಓದುತ್ತಿದ್ದೆ}}? (\textbf{\kn{ಓದುತ್ತ}} and \textbf{\kn{ಇದ್ದೆ}}.) What tells you \textquotedblleft{}was doing\textquotedblright{} from \textquotedblleft{}used to do\textquotedblright{}? (\textbf{The time word: one point in time, or }\kn{ಪ್ರತಿದಿನ}\textbf{.})
\end{tcolorbox}
\section[\texorpdfstring{\kn{ಓದುವಾಗ} (ōduvāga)}{ಓದುವಾಗ (ōduvāga)}]{\kn{ಓದುವಾಗ} (ōduvāga) — the scene, the event, and what was there}

\label{lesson:KA-C309-oduvaga}



\begin{quote}
Say \textquotedblleft{}when (somebody) comes\textquotedblright{}. (\emph{bandāga}.) Then \textquotedblleft{}I was reading\textquotedblright{}. (\emph{ōduttidde}.)

A told story often has one thing going on and another that cuts into it.
\end{quote}

\subsection*{You'll want to know: -\kn{ಉವಾಗ}, while it was going on}
\begin{quote}
\textbf{\kn{ನಾನು} \kn{ಊಟ} \kn{ಮಾಡುವಾಗ} \kn{ಫೋನ್} \kn{ಬಂತು}.}
\end{quote}

\begin{quote}
\emph{nānu ūṭa māḍuvāga phōn bantu.} — \textquotedblleft{}While I was having a meal, a call came.\textquotedblright{}
\end{quote}

\textbf{\kn{ಮಾಡುವಾಗ}} is \textbf{\kn{ಮಾಡುವ}}, \textquotedblleft{}that does\textquotedblright{}, with \textbf{\kn{ಆಗ}}, \textquotedblleft{}when\textquotedblright{}, the \textbf{\kn{ಆಗ}} you know from \textbf{\kn{ಬಂದಾಗ}}. \textbf{\kn{ಬಂದಾಗ}} is \textquotedblleft{}when (it) came\textquotedblright{}: done. \textbf{\kn{ಮಾಡುವಾಗ}} is \textquotedblleft{}while doing\textquotedblright{}: still going on.

That clause is the \textbf{scene}. The thing that cut into it is the \textbf{event}, in the plain past: \textbf{\kn{ಬಂತು}}. Kannada says a call \textquotedblleft{}came\textquotedblright{}, where English says the phone rang.

\begin{quote}
\textbf{\kn{ನಾವು} \kn{ನಡೆಯುವಾಗ} \kn{ಮಳೆ} \kn{ಬಂತು}.}
\end{quote}

\begin{quote}
\emph{nāvu naḍeyuvāga maḷe bantu.} — \textquotedblleft{}While we were walking, it rained.\textquotedblright{}
\end{quote}

A stem that ends in \textbf{-\kn{ಎ}} takes a \textbf{-\kn{ಯ}-} first: \textbf{\kn{ನಡೆ}} $\to$ \textbf{\kn{ನಡೆಯುವಾಗ}}. The \textbf{-\kn{ಉವಾಗ}} form carries no person: the rest of the sentence says who.

\subsection*{You'll want to know: \kn{ಇತ್ತು}, 'there was'}
\begin{quote}
\textbf{\kn{ನಮ್ಮ} \kn{ಊರಿನಲ್ಲಿ} \kn{ಒಂದು} \kn{ಚಿಕ್ಕ} \kn{ಅಂಗಡಿ} \kn{ಇತ್ತು}.}
\end{quote}

\begin{quote}
\emph{namma ūrinalli ondu cikka aṅgaḍi ittu.} — \textquotedblleft{}There was a small shop in our town.\textquotedblright{}
\end{quote}

\textbf{\kn{ನಮ್ಮ}} (\emph{namma}, \textquotedblleft{}our\textquotedblright{}); \textbf{\kn{ಊರಿನಲ್ಲಿ}} (\emph{ūrinalli}, \textquotedblleft{}in the town\textquotedblright{}). \textbf{\kn{ಇತ್ತು}} is the \textquotedblleft{}it\textquotedblright{} form of \textbf{\kn{ಇದ್ದೆ}}: \textquotedblleft{}there was\textquotedblright{}, \textquotedblleft{}it was\textquotedblright{}. Set it against the present:

\begin{quote}
\textbf{\kn{ಈಗ} \kn{ಅಲ್ಲಿ} \kn{ಒಂದು} \kn{ದೊಡ್ಡ} \kn{ಹೋಟೆಲ್} \kn{ಇದೆ}.}
\end{quote}

\begin{quote}
\emph{īga alli ondu doḍḍa hōṭel ide.} — \textquotedblleft{}Now there is a big hotel there.\textquotedblright{}
\end{quote}

\textquotedblleft{}There wasn't\textquotedblright{} is \textbf{\kn{ಇರಲಿಲ್ಲ}} (\emph{iralilla}), which ends in \textbf{\kn{ಇಲ್ಲ}}, \textquotedblleft{}not\textquotedblright{}. The same ending gives any verb's \textquotedblleft{}did not\textquotedblright{}, and it stays the same for every person: \textbf{\kn{ಅವನು} \kn{ಬರಲಿಲ್ಲ}} (\emph{avanu baralilla}, \textquotedblleft{}he did not come\textquotedblright{}), \textbf{\kn{ನಾನು} \kn{ಹೋಗಲಿಲ್ಲ}} (\emph{nānu hōgalilla}, \textquotedblleft{}I did not go\textquotedblright{}).

\subsection*{Your turn}
Say these aloud:

\begin{itemize}
\raggedright
  \item \textquotedblleft{}While I was having a meal, a call came\textquotedblright{} — \emph{nānu ūṭa māḍuvāga phōn bantu}
  \item \textquotedblleft{}While we were walking, it rained\textquotedblright{} — \emph{nāvu naḍeyuvāga maḷe bantu}
  \item \textquotedblleft{}There was a small shop in our town. Now there is a big hotel there.\textquotedblright{} — \emph{namma ūrinalli ondu cikka aṅgaḍi ittu. īga alli ondu doḍḍa hōṭel ide.}
  \item \textquotedblleft{}He did not come\textquotedblright{} — \emph{avanu baralilla}
\end{itemize}

\begin{tcolorbox}[breakable,colback=teal!4,colframe=teal!35!black,title={Before you move on}]
Which form paints the scene? (\textbf{-\kn{ಉವಾಗ}}: \textbf{\kn{ಮಾಡುವಾಗ}}, \textbf{\kn{ನಡೆಯುವಾಗ}}.) Which form drops in the event? (\textbf{The plain past}: \textbf{\kn{ಬಂತು}}.) \textquotedblleft{}There was\textquotedblright{}? (\textbf{\kn{ಇತ್ತು}}.) \textquotedblleft{}There wasn't\textquotedblright{}? (\textbf{\kn{ಇರಲಿಲ್ಲ}}.)
\end{tcolorbox}
\section[(review)]{Review — the past three ways}

\label{lesson:KA-R309-the-past-three-ways}



\begin{quote}
Say \textquotedblleft{}Yesterday I went to my home town.\textquotedblright{} (\emph{ninne nānu ūrige hōde}.) Then \textquotedblleft{}she gave\textquotedblright{} and \textquotedblleft{}it happened\textquotedblright{}. (\emph{koṭṭaḷu}, \emph{āyitu}.)
\end{quote}

\subsection*{Your turn}
Say these aloud:

\begin{itemize}
\raggedright
  \item We drank coffee, and afterwards we went to the cinema.
  \item While I was reading a book, my younger sister phoned.
  \item At nine last night it was raining.
  \item My grandmother used to tell a story every evening.
  \item There was a small shop in our village. Now there is a big hotel.
  \item one thing you were doing yesterday evening, with -\kn{ಉವಾಗ}, and one thing that cut into it
\end{itemize}

\begin{tcolorbox}[breakable,colback=teal!4,colframe=teal!35!black,title={Before you move on}]
Which past is the event? (\textbf{The plain past}: \textbf{\kn{ಹೋದೆ}}, \textbf{\kn{ಬಂತು}}.) Which is the scene? (\textbf{-\kn{ಉವಾಗ}}, or \textbf{-\kn{ಉತ್ತಿದ್ದೆ}} with one point in time.) Which is the habit? (\textbf{-\kn{ಉತ್ತಿದ್ದೆ}} with \textbf{\kn{ಪ್ರತಿದಿನ}}.) \textquotedblleft{}There was\textquotedblright{}? (\textbf{\kn{ಇತ್ತು}}.)
\end{tcolorbox}
